\section{The optional layer created by \emph{Cum sanctissima}}

The Congregation for the Doctrine of the Faith issued the decree \emph{Cum sanctissima} on 22 February 2020, effective 19 March 2020. Its purpose was to permit the celebration of saints canonized after the 1960 cut-off within celebrations according to the 1962 books. The accompanying official note stresses that this is an optional mechanism: a celebrant may continue to follow the sanctoral exactly as printed in the older books. It is therefore an overlay, not a revised edition of the 1962 general calendar.

\subsection{The eight operative rules}

\begin{enumerate}[leftmargin=1.7em,itemsep=0.45em]
\item For a just cause, a \emph{missa festiva latiore sensu} may be celebrated on class-III days other than the protected feasts in the decree's annex, and on a class-III vigil of a saint.
\item The Mass of any saint canonized after 26 July 1960 may be used on the date of that saint's universal liturgical memorial; a votive Mass may also be celebrated within the existing votive-Mass rubrics.
\item When that festal Mass is used, the complete corresponding Divine Office may be celebrated as an ordinary Office.
\item The omitted class-III feast or vigil is always commemorated, together with other occurring commemorations, subject to General Rubrics n. 111(d).
\item Formularies are selected in this order: the existing \emph{Proprium Sanctorum pro aliquibus locis}; a new supplement approved by the Holy See; otherwise the appropriate Common. A formulary supplied by the Proper or Supplement is used on the date assigned there.
\item On class-III or class-IV days, an optional ordinary commemoration may instead be made of a saint or mystery listed for that date in the Proper or approved Supplement, again subject to n. 111(d).
\item In houses of institutes of consecrated life or societies of apostolic life, the house superior---not the individual celebrant---decides how the provisions are used for the conventual Mass or Office in common or choir.
\item The seventy class-III feasts in the annex may not be impeded or omitted through these options. They may themselves be celebrated on class-III ferias of Lent or Passiontide, with the feria commemorated.
\end{enumerate}

The official note adds two especially important guardrails: class-I and class-II celebrations are unaffected, and the option is not a general power to invent a feast date or compose a proper. The decree works through the old rubrical categories and specified sources.

\subsection{The seventy protected class-III feasts}

This is the decree's complete annex. Titles below follow the Holy See's published English version; the Latin fixed tables above remain the edition baseline.

\begin{longtable}{@{}p{0.15\linewidth}p{0.77\linewidth}@{}}
\toprule
\textbf{Month} & \textbf{Protected dates and feasts}\\
\midrule
\endfirsthead
\toprule
\textbf{Month} & \textbf{Protected dates and feasts}\\
\midrule
\endhead
\bottomrule
\endfoot
January & 17, St. Anthony, Abbot; 20, Ss. Fabian and Sebastian; 21, St. Agnes; 24, St. Timothy; 25, Conversion of St. Paul; 26, St. Polycarp; 27, St. John Chrysostom; 29, St. Francis de Sales; 31, St. John Bosco.\\
February & 1, St. Ignatius; 5, St. Agatha; 6, St. Titus.\\
March & 6, Ss. Perpetua and Felicity; 7, St. Thomas Aquinas; 9, St. Frances of Rome; 12, St. Gregory I; 21, St. Benedict; 24, St. Gabriel the Archangel.\\
April & 11, St. Leo I; 14, St. Justin; 30, St. Catherine of Siena.\\
May & 2, St. Athanasius; 4, St. Monica; 5, St. Pius V; 9, St. Gregory of Nazianzus; 25, St. Gregory VII; 26, St. Philip Neri.\\
June & 5, St. Boniface; 11, St. Barnabas; 13, St. Anthony of Padua; 14, St. Basil the Great; 21, St. Aloysius Gonzaga; 30, Commemoration of St. Paul.\\
July & 7, Ss. Cyril and Methodius; 14, St. Bonaventure; 19, St. Vincent de Paul; 22, St. Mary Magdalene; 29, St. Martha; 31, St. Ignatius.\\
August & 2, St. Alphonsus Liguori; 4, St. Dominic; 5, Dedication of St. Mary of the Snows; 8, St. John Vianney; 12, St. Clare; 20, St. Bernard; 28, St. Augustine; 29, Beheading of St. John the Baptist.\\
September & 3, St. Pius X; 12, Most Holy Name of Mary; 16, Ss. Cornelius and Cyprian; 27, Ss. Cosmas and Damian; 30, St. Jerome.\\
October & 2, Holy Guardian Angels; 3, St. Thérèse of the Child Jesus; 4, St. Francis; 6, St. Bruno; 14, St. Callistus I; 15, St. Teresa.\\
November & 4, St. Charles; 11, St. Martin; 14, St. Josaphat; 18, Dedication of the Basilicas of Ss. Peter and Paul; 22, St. Cecilia; 23, St. Clement I; 24, St. John of the Cross.\\
December & 3, St. Francis Xavier; 6, St. Nicholas; 7, St. Ambrose; 11, St. Damasus I; 13, St. Lucy.\\
\end{longtable}

The annex contains exactly seventy entries. Protection here is protection against displacement by the decree's new options; it does not elevate the feasts to class II or override the ordinary precedence rules for unrelated occurrences.
